\documentclass[aspectratio=169,10pt]{beamer}
\usetheme{Madrid}
\usecolortheme{seagull}
\usepackage[utf8]{inputenc}
\usepackage{booktabs}
\usepackage{xcolor}
\usepackage{colortbl}
\setbeamertemplate{navigation symbols}{}
\setbeamertemplate{caption}{\raggedright\insertcaption\par}

% --- Speaker notes ---
% Default build = clean slides only.
% To export a notes handout, change the next line to:
%   \setbeameroption{show only notes}
% and recompile to a separate PDF.
\setbeamertemplate{note page}{\insertnote}

\definecolor{hl}{HTML}{1A6E5A}
\newcommand{\hl}[1]{\textcolor{hl}{\textbf{#1}}}

\title[Universality of Electricity Demand]{The Universality of Electricity Demand}
\subtitle{Zero-Shot Foundation Models Match Locally-Trained Accuracy on Brazil's Power Grid}
\author{Nelson Barlow}
\institute{Centro de Informática --- UFPE}
\date{}

\begin{document}

% ---------- 1. Title (~20s) ----------
\begin{frame}
  \titlepage
\end{frame}
\note{
Olá. Meu projeto investiga uma pergunta simples mas com grande impacto prático:
será que modelos de fundação de séries temporais conseguem prever a demanda de energia
do Brasil SEM nunca terem sido treinados em dados brasileiros? Eu chamo essa hipótese de
"universalidade da demanda elétrica". [~20s]
}

% ---------- 2. The problem (~30s) ----------
\begin{frame}{O problema: previsão de carga exige dados locais}
  \begin{itemize}
    \item Previsão de carga de curto prazo (STLF) sustenta despacho, reservas e o mercado day-ahead.
    \item Paradigma dominante: \hl{cada concessionária treina seu próprio modelo} com anos de histórico local + clima + calendário.
    \item Barreira pesada em redes com poucos dados: regiões recém-interligadas, microrredes, países em desenvolvimento.
  \end{itemize}
\end{frame}
\note{
Hoje, prever carga elétrica significa que cada operador constrói um modelo sob medida,
treinado em anos de dados locais. Isso funciona, mas é caro e impossível onde não há histórico:
uma microrrede nova, uma região recém-conectada. A pergunta é: dá pra pular essa etapa? [~30s]
}

% ---------- 3. Hypothesis (~40s) ----------
\begin{frame}{A hipótese: universalidade da demanda}
  \begin{columns}[T]
    \begin{column}{0.5\textwidth}
      \textbf{Forças universais} \\(transferem entre redes)
      \begin{itemize}
        \item Física de aquecimento/refrigeração
        \item Ritmo circadiano (pico diurno, vale noturno)
        \item Ciclo semanal trabalho/descanso
      \end{itemize}
    \end{column}
    \begin{column}{0.5\textwidth}
      \textbf{Forças locais} \\(NÃO transferem)
      \begin{itemize}
        \item Feriados nacionais
        \item Restrições hidrológicas
        \item Sazonalidade do hemisfério sul
      \end{itemize}
    \end{column}
  \end{columns}
  \vspace{1em}
  \begin{center}
    \hl{Se um modelo pré-treinado internalizou a parte universal, ele transfere zero-shot}\\
    e o erro residual fica confinado à parte local.
  \end{center}
\end{frame}
\note{
A intuição: a demanda elétrica tem duas camadas. Uma universal — física e comportamento humano,
que produz o mesmo formato de curva em qualquer lugar do mundo. E uma local — feriados, hidrologia,
estação invertida. Minha hipótese é que um modelo treinado em séries globais já aprendeu a camada
universal. Então ele deveria acertar zero-shot, e errar só na parte local. Vou testar exatamente isso. [~40s]
}

% ---------- 4. Setup (~35s) ----------
\begin{frame}{Dados e modelos}
  \begin{columns}[T]
    \begin{column}{0.52\textwidth}
      \textbf{Dados}
      \begin{itemize}
        \item ONS (dados.ons.org.br), CC-BY 4.0
        \item Horário, 2019--2025, 61.368 linhas
        \item 4 subsistemas do SIN (SE $\approx$ 55\%)
        \item Réplica: TEPCO (Tóquio)
        \item Teste: últimos 365 dias; zero-shot só vê o contexto
      \end{itemize}
    \end{column}
    \begin{column}{0.48\textwidth}
      \textbf{Modelos}
      \begin{itemize}
        \item \hl{Zero-shot:} Chronos-2 (120M), Moirai 2.0 (11M), TiRex (35M)
        \item \hl{Treinados:} N-BEATS (ajustado), Linear, LSTM
        \item Baseline: Naive (7 dias atrás)
      \end{itemize}
      \vspace{0.5em}
      \footnotesize Sem vazamento: ONS ausente do corpus do Chronos.
    \end{column}
  \end{columns}
\end{frame}
\note{
Uso dados horários da ONS, sete anos, os quatro subsistemas do sistema interligado. Comparo três
modelos de fundação zero-shot — Chronos-2, Moirai e TiRex — contra um N-BEATS ajustado e treinado
em cinco anos de dados locais, mais baselines. Importante: verifiquei que os dados da ONS NÃO estão
no corpus de treino do Chronos — não há vazamento. Os zero-shot só enxergam a janela de contexto. [~35s]
}

% ---------- 5. Main result (~50s) ----------
\begin{frame}{Resultado principal --- SE, horizonte 24h}
  \begin{center}
  \begin{tabular}{llr}
    \toprule
    Modelo & Tipo & MAPE \\
    \midrule
    Chronos-2 (fine-tuned) & Fine-tuned & 1{,}73\% \\
    \rowcolor{hl!15} \textbf{Chronos-2} & \textbf{Zero-shot} & \textbf{1{,}86\%} \\
    Moirai 2.0 & Zero-shot & 1{,}93\% \\
    N-BEATS (ajustado) & Treinado 5+ anos & 1{,}91\% $\pm$ 0{,}08\% \\
    TiRex & Zero-shot & 2{,}33\% \\
    Naive (7 dias) & Baseline & 5{,}13\% \\
    \bottomrule
  \end{tabular}
  \end{center}
  \vspace{0.5em}
  \begin{center}
    \hl{Zero-shot Chronos-2 (1,86\%) iguala um N-BEATS treinado em 5 anos (1,91\%)} --- sem treino local.
  \end{center}
\end{frame}
\note{
Este é o resultado central. No subsistema Sudeste, horizonte de 24 horas: o Chronos-2 zero-shot
atinge 1,86\% de MAPE. O N-BEATS, treinado em cinco anos de dados locais, fica em 1,91\%. Ou seja,
o modelo que NUNCA viu dados brasileiros empata com — na verdade, levemente supera — o modelo treinado
localmente. Todos os modelos de fundação batem o baseline naive de 5,13\% com folga. Fine-tuning leve
ganha só mais 7\%. [~50s]
}

% ---------- 6. Statistical tie (~30s) ----------
\begin{frame}{É um empate estatístico, não sorte}
  \begin{itemize}
    \item Teste de \hl{Diebold--Mariano} contra as 3 sementes do N-BEATS: \hl{$p > 0{,}29$} (sem diferença significativa).
    \item Intervalos de confiança 95\% por bootstrap se sobrepõem (Chronos-2: [1,70\%, 2,04\%]).
    \item N-BEATS foi devidamente ajustado: grid search (input length $\times$ lr), 3 sementes $\to$ 1,91\% $\pm$ 0,08\%.
  \end{itemize}
  \vspace{0.5em}
  \begin{center}
    Mesma acurácia. \hl{Um precisa de 5 anos de treino; o outro, de zero.}
  \end{center}
\end{frame}
\note{
Alguém poderia dizer: 1,86 contra 1,91, é ruído. Exatamente — e é esse o ponto. O teste de
Diebold-Mariano dá p maior que 0,29, então a diferença NÃO é estatisticamente significativa. Os
intervalos de confiança se sobrepõem. E garanti que o N-BEATS estava bem ajustado, com busca em
grade e três sementes. A conclusão honesta é: mesma acurácia. A diferença é que um exigiu cinco
anos de treino local e o outro, nenhum. [~30s]
}

% ---------- 7. Generalization (~45s) ----------
\begin{frame}{Generaliza em todas as dimensões}
  \begin{columns}[T]
    \begin{column}{0.5\textwidth}
      \begin{itemize}
        \item \hl{4 subsistemas:} Chronos-2 lidera todos (1,67--3,17\%), R$^2 > 0{,}90$.
        \item \hl{3 anos de teste:} estável, 1,89\% $\pm$ 0,04\%.
        \item \hl{Contexto:} limiar de 1 semana; MAPE cai pela metade ao ver um ciclo semanal completo. 30 dias = ótimo.
      \end{itemize}
    \end{column}
    \begin{column}{0.5\textwidth}
      \begin{itemize}
        \item \hl{Horizonte:} melhor em 24--168h; naive vence só após $\sim$2 semanas.
        \item \hl{Escala:} família Bolt escala log-linear (Tiny 2,28\% $\to$ Chronos-2 1,86\%).
        \item \hl{Probabilístico:} bem calibrado; Chronos-2 melhor CRPS, Moirai intervalos mais nítidos.
      \end{itemize}
    \end{column}
  \end{columns}
\end{frame}
\note{
O resultado não é um acaso do Sudeste. Ele se mantém nos quatro subsistemas, com R-quadrado acima de
0,90 em todos. Mantém-se em três anos diferentes de teste. Sobre contexto: o modelo precisa de pelo
menos uma semana de histórico — quando vê um ciclo semanal completo, o erro cai pela metade; 30 dias
é o ponto ideal. Em horizonte, os modelos de fundação ganham até duas semanas; depois disso, o naive
volta a vencer. Use para operação de curto prazo. [~45s]
}

% ---------- 8. Cross-country (~40s) ----------
\begin{frame}{Réplica entre países: Tóquio (TEPCO)}
  \begin{center}
  \begin{tabular}{llr}
    \toprule
    Modelo & Tipo & MAPE \\
    \midrule
    \rowcolor{hl!15} Chronos-2 & Zero-shot & \textbf{3{,}91\%} \\
    Moirai 2.0 & Zero-shot & 3{,}94\% \\
    TiRex & Zero-shot & 4{,}05\% \\
    N-BEATS & Treinado 5 anos & 4{,}44\% \\
    Naive (7 dias) & Baseline & 8{,}86\% \\
    \bottomrule
  \end{tabular}
  \end{center}
  \vspace{0.4em}
  \begin{center}
    Outro país, \hl{outro hemisfério, outro calendário} --- mesmo ranking.\\
    Os três zero-shot \hl{vencem} o N-BEATS treinado localmente.
  \end{center}
\end{frame}
\note{
A evidência mais forte da universalidade: repeti tudo na rede de Tóquio, a TEPCO. País diferente,
hemisfério diferente, calendário de feriados diferente. E o mesmo padrão aparece: os três modelos
zero-shot batem o N-BEATS treinado em cinco anos de dados japoneses. O MAPE absoluto é maior porque
a demanda japonesa é mais volátil, mas o resultado qualitativo se sustenta entre países. Isso é difícil
de explicar a não ser que exista mesmo uma estrutura universal na demanda. [~40s]
}

% ---------- 9. When it fails (~40s) ----------
\begin{frame}{Onde falha: feriados}
  \begin{itemize}
    \item Em feriados nacionais brasileiros, o Chronos-2 degrada para \hl{4--5$\times$} o erro de dia normal.
    \item Motivo: nada no corpus global sinaliza um feriado brasileiro $\to$ ele prevê demanda de dia útil.
    \item É propriedade \hl{dos dados, não de um modelo}: Moirai 7,9\%, TiRex 9,1\%, naive 12,8\% nos mesmos feriados (vs. $\sim$1,6\% normal).
  \end{itemize}
  \vspace{0.5em}
  \begin{center}
    Exatamente o previsto: \hl{a fronteira da universalidade é o calendário local.}
  \end{center}
\end{frame}
\note{
E onde o modelo falha? Exatamente onde a hipótese previu: feriados. Em feriados nacionais, o erro
do Chronos-2 sobe de cerca de 1,6\% para quatro ou cinco vezes isso. A razão é intuitiva — nada no
corpus global de treino diz a ele que é feriado no Brasil, então ele prevê um dia útil comum. E não é
defeito de um modelo só: todos degradam igual, inclusive o naive. Isso confirma a fronteira: o que não
transfere é o conhecimento de calendário local. [~40s]
}

% ---------- 10. The fix (~35s) ----------
\begin{frame}{A correção funciona: 1 covariável de feriado}
  \begin{center}
  \begin{tabular}{lrrr}
    \toprule
    Regime & Zero-shot & + Covariável & $\Delta$ \\
    \midrule
    \rowcolor{hl!15} Feriado & 6{,}76\% & \textbf{4{,}37\%} & \hl{$-$35\%} \\
    Carnaval & 5{,}68\% & 4{,}64\% & $-$18\% \\
    Dia seguinte & 2{,}29\% & 1{,}93\% & $-$16\% \\
    Dia normal & 1{,}53\% & 1{,}53\% & 0\% \\
    Geral & 1{,}82\% & 1{,}73\% & $-$5\% \\
    \bottomrule
  \end{tabular}
  \end{center}
  \vspace{0.4em}
  \begin{center}
    Uma covariável binária de feriado + fine-tuning leve \hl{recupera 35\% do erro de feriado},\\
    sem tocar nos dias normais.
  \end{center}
\end{frame}
\note{
E se a falha é falta de conhecimento de calendário, basta fornecer esse conhecimento. Adicionei UMA
covariável binária de feriado, com fine-tuning leve. O erro de feriado cai 35\%, o Carnaval 18\%, o
dia seguinte ao feriado 16\% — e os dias normais ficam intactos. Isso fecha o argumento: a fronteira
da universalidade é conhecimento de calendário, não uma limitação de arquitetura. Forneça o pedaço
local que falta, e o modelo recupera quase toda a diferença. [~35s]
}

% ---------- 11. Takeaway (~35s) ----------
\begin{frame}{Conclusão e receita operacional}
  \begin{block}{Universalidade limitada da demanda elétrica}
    Estrutura física e comportamental transfere zero-shot entre redes, anos e países;
    só o calendário local não transfere --- e uma covariável o recupera.
  \end{block}
  \textbf{Receita para redes com poucos dados:}
  \begin{enumerate}
    \item Comece zero-shot (Chronos-2, ou Moirai 11M se houver restrição) com $\sim$30 dias de contexto.
    \item Fine-tune se possível ($\sim$7\% de ganho, 40 min em CPU).
    \item Adicione covariável binária de feriado.
    \item Não use além de $\sim$2 semanas de horizonte.
  \end{enumerate}
  \vspace{0.3em}
  \footnotesize Código e dados: \texttt{github.com/nelsonbarlow/brazil-energy-forecast}
\end{frame}
\note{
Em resumo: a demanda elétrica tem uma universalidade, mas com fronteira precisa. A estrutura física
e comportamental transfere entre redes, anos e até países. Só o calendário local não transfere — e uma
única covariável o recupera. Para um operador com poucos dados, a receita prática é: comece zero-shot
com 30 dias de contexto, faça fine-tuning se puder, adicione a covariável de feriado, e não passe de
duas semanas de horizonte. Previsão day-ahead precisa não exige mais anos de desenvolvimento local.
Obrigado. Todo o código está no GitHub. [~35s]
}

\end{document}
